\chapter{Grensorbitalen en reactiviteit}\label{ch:b2:frontier-orbitals}

Een fles cyclopentadieen die op de plank blijft staan, verandert langzaam in haar dimeer:
twee moleculen van hetzelfde \omterm{def:b2:frontier-orbitals:diels-alder}{dieen} verbinden zich tot een overbrugde bicyclische verbinding, en
er ontstaat maar één van de mogelijke stereo-isomeren. Niets in de gebogen pijlen van
het deel van jaar~1 voorspelt welke. Twee orbitalen doen dat wel: de hoogste bezette
orbitaal van het ene molecuul en de laagste lege orbitaal van het andere. Dit
hoofdstuk maakt van de wisselwerking tussen twee niveaus uit \cref{ch:b2:diatomic-mos} een
werktuig voor reactiviteit --- welk atoom een nucleofiel aanvalt, van welke kant, en
waarom een \omterm{def:b2:frontier-orbitals:diels-alder}{dieen} en een alkeen in één stap tot een ring combineren.

\begin{recall}
\Cref{ch:b2:diatomic-mos}: twee wisselwerkende orbitalen, afstoting met vier elektronen.
\Cref{ch:b2:huckel}: hückelniveaus en -coëfficiënten van geconjugeerde systemen.
Het deel van jaar~1: nucleofielen en elektrofielen, gebogen pijlen, kinetische
controle, stereospecifieke, stereoselectieve en regioselectieve reacties,
elektrofiele addities aan alkenen, nucleofiele substitutie.
\end{recall}

\section{Twee orbitalen, opnieuw bekeken}

Wanneer twee moleculen elkaar naderen, overlappen hun orbitalen en wisselwerken ze. Het grootste deel van
de wisselwerking is zwak, omdat de orbitalen ver uit elkaar liggen in energie; een storingsvorm van
het resultaat voor twee niveaus volstaat dan.

\begin{theorem}[Zwakke wisselwerking van twee niveaus]\label{thm:b2:frontier-orbitals:perturbation}
Twee orbitalen met energieën $\varepsilon_1 < \varepsilon_2$, gekoppeld door een \omterm{def:b2:diatomic-mos:integrals}{resonantie-integraal}
$\beta$ met $|\beta| \ll \Delta\varepsilon = \varepsilon_2 - \varepsilon_1$ (overlap verwaarloosd), worden
verschoven naar
\[
  \varepsilon_1 - \frac{\beta^2}{\Delta\varepsilon}, \qquad \varepsilon_2 + \frac{\beta^2}{\Delta\varepsilon}
\]
tot op eerste orde in $\beta^2/\Delta\varepsilon^2$: het lagere niveau daalt en het hogere stijgt met hetzelfde
bedrag, des te kleiner naarmate de afstand groter is.
\end{theorem}

\begin{proof}
Volgens \cref{thm:b2:diatomic-mos:heteronuclear} zijn de exacte niveaus $\bar\varepsilon \mp
\sqrt{\Delta\varepsilon^2/4 + \beta^2}$. Met $u = 4\beta^2/\Delta\varepsilon^2 \ll 1$ is $\sqrt{\Delta\varepsilon^2/4
+ \beta^2} = \frac{\Delta\varepsilon}{2}\sqrt{1 + u} \approx \frac{\Delta\varepsilon}{2}(1 + u/2) =
\frac{\Delta\varepsilon}{2} + \frac{\beta^2}{\Delta\varepsilon}$, en $\bar\varepsilon \mp \Delta\varepsilon/2 = \varepsilon_1,
\varepsilon_2$.
\end{proof}

\begin{proposition}[Twee elektronen trekken aan, vier stoten af]\label{prop:b2:frontier-orbitals:four-electron}
Als de twee wisselwerkende orbitalen samen twee elektronen bevatten, wordt het systeem gestabiliseerd
met ongeveer $2\beta^2/\Delta\varepsilon$. Als ze er vier bevatten, wordt het systeem gedestabiliseerd.
\end{proposition}

\begin{proof}
Twee elektronen gaan naar het verlaagde niveau: winst $2\beta^2/\Delta\varepsilon$. Met vier wordt ook het hogere niveau
gevuld; tot op de orde van \cref{thm:b2:frontier-orbitals:perturbation} heffen de twee
verschuivingen elkaar op, en met de overlap meegerekend stijgt het hogere niveau meer dan het lagere
daalt (\cref{prop:b2:diatomic-mos:asymmetry}): een netto afstoting.
\end{proof}

\begin{omfigure}
\begin{tikzpicture}[font=\small]
  \begin{scope}
    \pic (a) at (0,0) {omlevel={ud}{0.7}};
    \pic (b) at (3,1.4) {omlevel={empty}{0.7}};
    \pic (lo) at (1.5,-0.5) {omlevel={ud}{0.7}};
    \pic (hi) at (1.5,1.9) {omlevel={empty}{0.7}};
    \draw[omcorr, densely dashed] (a-r) -- (lo-l) (a-r) -- (hi-l) (b-l) -- (lo-r) (b-l) -- (hi-r);
    \node[below=6pt] at (1.5,-0.6) {twee elektronen: gestabiliseerd};
    \node[left] at (a-l) {gevuld};
    \node[right] at (b-r) {leeg};
  \end{scope}
  \begin{scope}[xshift=7cm]
    \pic (a) at (0,0) {omlevel={ud}{0.7}};
    \pic (b) at (3,0.6) {omlevel={ud}{0.7}};
    \pic (lo) at (1.5,-0.5) {omlevel={ud}{0.7}};
    \pic (hi) at (1.5,1.25) {omlevel={ud}{0.7}};
    \draw[omcorr, densely dashed] (a-r) -- (lo-l) (a-r) -- (hi-l) (b-l) -- (lo-r) (b-l) -- (hi-r);
    \node[below=6pt] at (1.5,-0.6) {vier elektronen: afgestoten};
    \node[left] at (a-l) {gevuld};
    \node[right] at (b-r) {gevuld};
  \end{scope}
\end{tikzpicture}

{\small Wisselwerking van twee orbitalen van naburige moleculen. Een gevulde orbitaal
tegenover een lege geeft een netto stabilisatie; twee gevulde orbitalen stoten elkaar af, waarbij het
hogere niveau meer stijgt dan het lagere daalt. Deze afstoting is de oorsprong
van de sterische hindering tussen moleculen.}
\end{omfigure}

\section{Grensorbitalen}

\begin{definition}[Grensorbitalen]\label{def:b2:frontier-orbitals:frontier}
De \emph{HOMO}\index{HOMO} (hoogste bezette \omterm{def:b2:diatomic-mos:lcao}{molecuulorbitaal}) en de
\emph{LUMO}\index{LUMO} (laagste onbezette \omterm{def:b2:diatomic-mos:lcao}{molecuulorbitaal}) van een molecuul zijn zijn
\emph{grensorbitalen}\index{grensorbitalen}.
\end{definition}

\begin{proposition}[De overheersende wisselwerking]\label{prop:b2:frontier-orbitals:dominant}
Tussen twee moleculen met gesloten schil zijn de stabiliserende wisselwerkingen die het meest tellen,
die tussen de \omterm{def:b2:frontier-orbitals:frontier}{HOMO} van het ene en de \omterm{def:b2:frontier-orbitals:frontier}{LUMO} van het andere; van de twee zulke paren overheerst het paar
met de kleinste afstand. Een nucleofiel reageert via zijn \omterm{def:b2:frontier-orbitals:frontier}{HOMO}, een
elektrofiel via zijn \omterm{def:b2:frontier-orbitals:frontier}{LUMO}.
\end{proposition}

\begin{proof}[Redenering]
Paren gevuld--gevuld stoten af (\cref{prop:b2:frontier-orbitals:four-electron}) en paren leeg--leeg
bevatten geen elektronen: alleen paren gevuld--leeg stabiliseren, elk met $2\beta^2/\Delta\varepsilon$.
De \omterm{def:b2:frontier-orbitals:frontier}{HOMO} is het hoogste gevulde niveau en de \omterm{def:b2:frontier-orbitals:frontier}{LUMO} het laagste lege, dus een
paar HOMO--LUMO heeft daaronder de kleinste afstand; volgens de stelling geeft de kleinste afstand
de grootste term.
\end{proof}

\begin{definition}[Ladingscontrole en orbitaalcontrole]\label{def:b2:frontier-orbitals:control}
Een reactie staat onder \emph{ladingscontrole}\index{ladingscontrole} wanneer de aantrekking
tussen de ladingen van de twee partners beslist waar ze reageren, en onder
\emph{orbitaalcontrole}\index{orbitaalcontrole} wanneer de overlap van de \omterm{def:b2:frontier-orbitals:frontier}{grensorbitalen}
dat beslist, waarbij de aanval plaatsvindt waar hun coëfficiënten het grootst zijn.
\end{definition}

\begin{definition}[Harde en zachte deeltjes]\label{def:b2:frontier-orbitals:hard-soft}
Een \emph{hard nucleofiel}\index{hard nucleofiel} of \emph{hard elektrofiel}
\index{hard elektrofiel} is klein, draagt een geconcentreerde lading en heeft
een \omterm{def:b2:frontier-orbitals:frontier}{HOMO} (respectievelijk een \omterm{def:b2:frontier-orbitals:frontier}{LUMO}) die ver ligt van die van zijn partners: het reageert onder \omterm{def:b2:frontier-orbitals:control}{ladingscontrole}.
Een \emph{zacht nucleofiel}\index{zacht nucleofiel} of \emph{zacht elektrofiel}
\index{zacht elektrofiel} is groot en polariseerbaar, met een hoge \omterm{def:b2:frontier-orbitals:frontier}{HOMO}
(respectievelijk een lage \omterm{def:b2:frontier-orbitals:frontier}{LUMO}): het reageert onder \omterm{def:b2:frontier-orbitals:control}{orbitaalcontrole}. Hard reageert bij voorkeur
met hard, zacht met zacht.
\end{definition}

Het hydroxide-ion en het fluoride-ion zijn \omterm{def:b2:frontier-orbitals:hard-soft}{harde nucleofielen}, jodide en thiolaten
zachte; een proton en een carbokation zijn \omterm{def:b2:frontier-orbitals:hard-soft}{harde elektrofielen}, het koolstofatoom van een
halogeenalkaan is een \omterm{def:b2:frontier-orbitals:hard-soft}{zacht elektrofiel}.

\begin{definition}[Ambident nucleofiel]\label{def:b2:frontier-orbitals:ambident}
Een \emph{ambident nucleofiel}\index{ambident nucleofiel} heeft twee nucleofiele atomen
die door conjugatie verbonden zijn en die elk een elektrofiel kunnen aanvallen.
\end{definition}

Het cyanide-ion wordt door \omterm{def:b2:frontier-orbitals:hard-soft}{harde elektrofielen} aangevallen op stikstof, waar de negatieve lading groter is,
en door zachte op koolstof, waar zijn \omterm{def:b2:frontier-orbitals:frontier}{HOMO} de grootste coëfficiënt heeft:
een halogeenalkaan geeft een nitril \ce{R-C#N}. Enolaationen, die in
\cref{ch:b2:enolates-aldol} aan bod komen, reageren op dezelfde manier op koolstof of op zuurstof.

\begin{method}[Een analyse met grensorbitalen]\label{met:b2:frontier-orbitals:fmo}
\begin{enumerate}
  \item Identificeer het nucleofiel (hoge \omterm{def:b2:frontier-orbitals:frontier}{HOMO}) en het elektrofiel (lage \omterm{def:b2:frontier-orbitals:frontier}{LUMO}).
  \item Teken de twee \omterm{def:b2:frontier-orbitals:frontier}{grensorbitalen} met hun coëfficiënten en tekens.
  \item De reactie vindt plaats waar de overlap het grootst is: atomen met de grootste
        coëfficiënten, lobben van hetzelfde teken tegenover elkaar.
  \item De naderingsrichting volgt de vorm van de \omterm{def:b2:frontier-orbitals:frontier}{LUMO} (loodrecht op een
        \ce{C=O}-vlak, aan de achterkant van een \ce{C-X}-binding).
  \item Een kleinere afstand HOMO--LUMO betekent een snellere reactie (als al het andere gelijk is).
\end{enumerate}
\end{method}

De richting van de aanval volgt uit de \omterm{def:b2:frontier-orbitals:frontier}{LUMO}. De \omterm{def:b2:frontier-orbitals:frontier}{LUMO} van een carbonylgroep is haar
$\pi^*$-orbitaal, groter op koolstof en loodrecht op het vlak: nucleofielen naderen
het koolstofatoom van boven of van onder het vlak, niet langs de as \ce{C=O}. De \omterm{def:b2:frontier-orbitals:frontier}{LUMO} van een
halogeenalkaan is de $\sigma^*$-orbitaal van de binding \ce{C-X}, waarvan de grote lob op koolstof
van X weg wijst: het nucleofiel valt aan langs de achterkant, en het koolstofatoom wordt geïnverteerd
--- de stereochemie van de bimoleculaire substitutie van het deel van jaar~1.

\section{De diels-alderreactie}

\begin{definition}[Gecoördineerde reactie]\label{def:b2:frontier-orbitals:concerted}
Een \emph{gecoördineerde reactie}\index{gecoordineerde reactie@gecoördineerde reactie} vormt en verbreekt al haar bindingen in
één enkele stap, via één overgangstoestand, zonder intermediair.
\end{definition}

\begin{definition}[Diels-alderreactie]\label{def:b2:frontier-orbitals:diels-alder}
De \emph{diels-alderreactie}\index{diels-alderreactie} is de gecoördineerde additie
van een geconjugeerd \emph{dieen}\index{dieen}, in zijn s-cisconformatie, aan een alkeen of
alkyn, het \emph{dienofiel}\index{dienofiel}, waarbij een zesring ontstaat met twee
nieuwe $\sigma$-bindingen en één nieuwe $\pi$-binding.
\end{definition}

\begin{omfigure}
\begin{tikzpicture}[font=\small, >={Stealth}, scale=1.25]
  % diene s-cis: C1 (0,1.6) C2 (0.6,2.4) C3 (1.6,2.4) C4 (2.2,1.6); dienophile below
  \draw[thick] (0,1.6) -- (0.6,2.4) -- (1.6,2.4) -- (2.2,1.6);
  \draw[thick] (0.2,1.62) -- (0.65,2.22);
  \draw[thick] (2.0,1.62) -- (1.55,2.22);
  \draw[thick] (0.2,0) -- (2.0,0); \draw[thick] (0.35,0.14) -- (1.85,0.14);
  \node[font=\scriptsize, anchor=east] at (-0.05,1.6) {1};
  \node[font=\scriptsize, anchor=south east] at (0.5,2.42) {2};
  \node[font=\scriptsize, anchor=south] at (1.6,2.42) {3};
  \node[font=\scriptsize, anchor=west] at (2.25,1.6) {4};
  % arrows: dienophile pi -> C1..C(dienophile); C1=C2 pi -> C2-C3; C3=C4 pi -> C4-C
  \draw[->, thick, omProp] (1.1,0.22) .. controls (0.9,0.9) and (0.3,0.9) .. (0.12,1.42);
  \draw[->, thick, omProp] (0.45,1.95) .. controls (0.6,2.75) and (1.0,2.8) .. (1.1,2.48);
  \draw[->, thick, omProp] (1.78,1.92) .. controls (2.5,1.6) and (2.4,0.6) .. (2.05,0.2);
  \node at (3.3,1.2) {$\longrightarrow$};
  \begin{scope}[xshift=4.2cm]
    \draw[thick] (0,1.6) -- (0.6,2.4) -- (1.6,2.4) -- (2.2,1.6) -- (2.0,0) -- (0.2,0) -- cycle;
    \draw[thick] (0.72,2.26) -- (1.48,2.26);
  \end{scope}
\end{tikzpicture}

{\small De \omterm{def:b2:frontier-orbitals:diels-alder}{diels-alderreactie} van butadieen en etheen, geschreven met gebogen pijlen:
drie elektronenparen verplaatsen zich tegelijk, in een cyclische, gecoördineerde stap; twee $\sigma$-bindingen
ontstaan aan de uiteinden van het \omterm{def:b2:frontier-orbitals:diels-alder}{dieen}, en de $\pi$-binding verschuift naar het midden ervan.}
\end{omfigure}

De overheersende wisselwerking is die tussen de \omterm{def:b2:frontier-orbitals:frontier}{HOMO} van het \omterm{def:b2:frontier-orbitals:diels-alder}{dieen}, $\psi_2$, en de \omterm{def:b2:frontier-orbitals:frontier}{LUMO} van het
\omterm{def:b2:frontier-orbitals:diels-alder}{dienofiel}, $\pi^*$. Hun lobben aan de uiteinden van het \omterm{def:b2:frontier-orbitals:diels-alder}{dieen} en op de twee koolstofatomen van het
\omterm{def:b2:frontier-orbitals:diels-alder}{dienofiel} hebben overeenstemmende tekens wanneer het \omterm{def:b2:frontier-orbitals:diels-alder}{dienofiel} het vlak van het \omterm{def:b2:frontier-orbitals:diels-alder}{dieen} nadert:
beide nieuwe bindingen kunnen tegelijk ontstaan, aan dezelfde kant van elke partner.

\begin{omfigure}
\begin{tikzpicture}[font=\small]
  \foreach \x/\c in {0/0.602, 1/0.372, 2/-0.372, 3/-0.602} {
    \pic at (\x,2) {omp={1.5*\c}{90}};
  }
  \draw[black!45] (-0.2,2) -- (3.2,2);
  \node[anchor=west] at (3.6,2) {HOMO van het dieen ($\psi_2$)};
  \pic at (0,0) {omp={-0.85}{90}}; \pic at (3,0) {omp={0.85}{90}};
  \draw[black!45] (-0.2,0) -- (3.2,0);
  \node[anchor=west] at (3.6,0) {LUMO van het dienofiel ($\pi^*$)};
  \draw[omProp, thick, densely dashed] (0,1.45) -- (0,0.6);
  \draw[omProp, thick, densely dashed] (3,1.45) -- (3,0.6);
  \node[font=\scriptsize, omProp, anchor=east] at (-0.15,1.0) {in fase};
  \node[font=\scriptsize, omProp, anchor=west] at (3.15,1.0) {in fase};
\end{tikzpicture}

{\small \omterm{def:b2:frontier-orbitals:frontier}{Grensorbitalen} van butadieen en etheen, met lobben in verhouding tot hun
hückelcoëfficiënten. Tegenover elkaar hebben de eindlobben van de \omterm{def:b2:frontier-orbitals:frontier}{HOMO} van het \omterm{def:b2:frontier-orbitals:diels-alder}{dieen} en
de lobben van de \omterm{def:b2:frontier-orbitals:frontier}{LUMO} van het \omterm{def:b2:frontier-orbitals:diels-alder}{dienofiel} aan beide uiteinden hetzelfde teken: de twee
$\sigma$-bindingen kunnen samen ontstaan, aan dezelfde kant van elke partner.}
\end{omfigure}

\begin{proposition}[Stereospecificiteit]\label{prop:b2:frontier-orbitals:da-stereospecific}
De \omterm{def:b2:frontier-orbitals:diels-alder}{diels-alderreactie} is stereospecifiek: substituenten die cis staan op het \omterm{def:b2:frontier-orbitals:diels-alder}{dienofiel}, staan
cis in het product, trans-substituenten blijven trans.
\end{proposition}

\begin{proof}
De reactie is gecoördineerd en beide nieuwe bindingen ontstaan aan dezelfde kant van het
\omterm{def:b2:frontier-orbitals:diels-alder}{dienofiel}: zijn twee koolstofatomen draaien nooit ten opzichte van elkaar, en de relatieve
posities van hun substituenten worden in de ring overgenomen.
\end{proof}

\begin{proposition}[Regioselectiviteit]\label{prop:b2:frontier-orbitals:da-regio}
Met asymmetrische partners verbindt het hoofdproduct het atoom van het \omterm{def:b2:frontier-orbitals:diels-alder}{dieen} met de
grootste coëfficiënt in zijn \omterm{def:b2:frontier-orbitals:frontier}{HOMO} met het atoom van het \omterm{def:b2:frontier-orbitals:diels-alder}{dienofiel} met de grootste
coëfficiënt in zijn \omterm{def:b2:frontier-orbitals:frontier}{LUMO}.
\end{proposition}

\begin{proof}[Redenering]
In de overgangstoestand zijn de twee nieuwe bindingen nog niet even ver gevormd; de
stabilisatie $2\beta^2/\Delta\varepsilon$ is groter wanneer de grotere coëfficiënten overlappen ($\beta$
groeit met het product van de coëfficiënten van de overlappende atomen), zodat die
oriëntatie bij een lagere energie bereikt wordt.
\end{proof}

\begin{omfigure}
\begin{tikzpicture}
\begin{axis}[ybar, width=0.47\linewidth, height=4.6cm, ymin=-0.8, ymax=0.8,
  xtick={0,1,2,3,4}, xticklabels={O,C1,C2,C3,C4}, axis x line=bottom, axis y line=left,
  extra y ticks={0}, extra y tick style={grid=major, grid style={black!50}}, enlarge x limits=0.15,
  tick label style={font=\scriptsize}, ylabel={coëfficiënt}, label style={font=\small},
  title={\small HOMO van 1-methoxybutadieen}, bar width=8pt]
\addplot[fill=omDef!50, draw=omDef] table[x=atom, y=c] {figdata/out/bachelor-2-huckel-c-diene.dat};
\end{axis}
\end{tikzpicture}\hfill
\begin{tikzpicture}
\begin{axis}[ybar, width=0.47\linewidth, height=4.6cm, ymin=-0.8, ymax=0.8,
  xtick={1,2,3,4}, xticklabels={C$\beta$,C$\alpha$,C(O),O}, axis x line=bottom, axis y line=left,
  extra y ticks={0}, extra y tick style={grid=major, grid style={black!50}}, enlarge x limits=0.15,
  tick label style={font=\scriptsize}, ylabel={coëfficiënt}, label style={font=\small},
  title={\small LUMO van propenal}, bar width=8pt]
\addplot[fill=omThm!45, draw=omThm] table[x=atom, y=c] {figdata/out/bachelor-2-huckel-c-dienophile.dat};
\end{axis}
\end{tikzpicture}

{\small Hückelmodel met parameters voor heteroatomen ($\alpha_{\ce{O}} = \alpha + 2\beta$,
$\beta_{\ce{CO}} = 0.8\beta$ voor de etherzuurstof; $\alpha_{\ce{O}} = \alpha + \beta$, $\beta_{\ce{CO}} =
\beta$ voor de carbonylzuurstof; een model). De \omterm{def:b2:frontier-orbitals:frontier}{HOMO} van het \omterm{def:b2:frontier-orbitals:diels-alder}{dieen} is het grootst op C4, het uiteinde
ver van de methoxygroep (0.60 tegenover 0.51); de \omterm{def:b2:frontier-orbitals:frontier}{LUMO} van het \omterm{def:b2:frontier-orbitals:diels-alder}{dienofiel} is het grootst op
het $\beta$-koolstofatoom (0.66). Die twee binden aan elkaar: de methoxy- en de aldehydegroep komen
op naburige koolstofatomen van de ring terecht.}
\end{omfigure}

De donorgroep verhoogt de \omterm{def:b2:frontier-orbitals:frontier}{HOMO} van het \omterm{def:b2:frontier-orbitals:diels-alder}{dieen} ($m = 0.48$ in plaats van 0.62) en de acceptor
verlaagt de \omterm{def:b2:frontier-orbitals:frontier}{LUMO} van het \omterm{def:b2:frontier-orbitals:diels-alder}{dienofiel} ($m = -0.35$ in plaats van $-1$): de afstand wordt kleiner, en zulke
paren reageren veel sneller dan butadieen met etheen. Een elektronenrijk \omterm{def:b2:frontier-orbitals:diels-alder}{dieen} en een
elektronenarm \omterm{def:b2:frontier-orbitals:diels-alder}{dienofiel} vormen de typische combinatie.

\begin{definition}[Endo- en exo-adducten]\label{def:b2:frontier-orbitals:endo}
Wanneer een cyclisch \omterm{def:b2:frontier-orbitals:diels-alder}{dieen} addeert aan een \omterm{def:b2:frontier-orbitals:diels-alder}{dienofiel} dat een onverzadigde substituent draagt, heeft het
\emph{endo-adduct}\index{endo-adduct} die substituent gericht naar de nieuw
gevormde dubbele binding (onder het \omterm{def:b2:frontier-orbitals:diels-alder}{dieen} in de overgangstoestand), het
\emph{exo-adduct}\index{exo-adduct} ervan weg. De \emph{endoregel}\index{endoregel} zegt
dat het endo-adduct sneller gevormd wordt.
\end{definition}

\begin{proposition}[De endoregel]\label{prop:b2:frontier-orbitals:endo}
Onder kinetische controle is het \omterm{def:b2:frontier-orbitals:endo}{endo-adduct} het hoofdproduct, hoewel het \omterm{def:b2:frontier-orbitals:endo}{exo-adduct}
vaak het stabielere is.
\end{proposition}

\admitted

De gebruikelijke verklaring is een secundaire overlap, bij de endonadering, tussen de
lobben van het $\pi$-systeem van de substituent en die van de centrale atomen van de \omterm{def:b2:frontier-orbitals:frontier}{HOMO} van
het \omterm{def:b2:frontier-orbitals:diels-alder}{dieen}; ze verlaagt de endo-overgangstoestand zonder een binding te vormen. De regel is
empirisch en kent uitzonderingen; het deel van jaar~3 behandelt deze reacties met de
symmetrie van hun orbitalen.

\begin{omfigure}
\begin{tikzpicture}[font=\small]
  \foreach \xs/\mode in {0/endo, 5/exo} {
    \begin{scope}[xshift=\xs cm]
      % cyclopentadiene seen from the side: C1, C2=C3, C4, CH2 bridge on top
      \draw[thick] (0,2) -- (0.6,2.5) -- (1.8,2.5) -- (2.4,2);
      \draw[thick] (0.75,2.38) -- (1.65,2.38);
      \draw[thick] (0,2) -- (1.2,3.2) -- (2.4,2);
      \node[font=\scriptsize, anchor=south] at (1.2,3.2) {\ce{CH2}};
      % dienophile
      \draw[thick] (0.2,0) -- (2.2,0); \draw[thick] (0.35,0.1) -- (2.05,0.1);
      \draw[omProp, densely dashed, thick] (0,1.9) -- (0.2,0.15);
      \draw[omProp, densely dashed, thick] (2.4,1.9) -- (2.2,0.15);
      \node[anchor=north] at (1.2,-1.3) {\mode};
    \end{scope}
  }
  % endo: carbonyls tucked under the diene
  \draw[thick] (0.2,0) -- (0.75,0.95) (2.2,0) -- (1.65,0.95);
  \node[font=\scriptsize, fill=white, inner sep=1pt] at (0.75,1.1) {C=O};
  \node[font=\scriptsize, fill=white, inner sep=1pt] at (1.65,1.1) {C=O};
  \draw[omThm, densely dotted, thick] (0.75,1.25) -- (0.65,2.35);
  \draw[omThm, densely dotted, thick] (1.65,1.25) -- (1.75,2.35);
  % exo: carbonyls pointing away
  \draw[thick] (5.2,0) -- (4.95,-0.85) (7.2,0) -- (7.45,-0.85);
  \node[font=\scriptsize] at (4.95,-1.0) {C=O};
  \node[font=\scriptsize] at (7.45,-1.0) {C=O};
\end{tikzpicture}

{\small Endo- en exonadering van maleïnezuuranhydride aan cyclopentadieen, schematisch.
De bindingen die ontstaan, zijn gestreept. Bij de endonadering liggen de carbonylgroepen onder het
\omterm{def:b2:frontier-orbitals:diels-alder}{dieen}, en hun $\pi$-orbitalen overlappen met de centrale koolstofatomen ervan (gestippeld, secundaire
overlap); bij de exonadering wijzen ze ervan weg.}
\end{omfigure}

\begin{method}[Een diels-alderproduct tekenen]\label{met:b2:frontier-orbitals:diels-alder}
\begin{enumerate}
  \item Zet het \omterm{def:b2:frontier-orbitals:diels-alder}{dieen} in zijn s-cisconformatie, met het \omterm{def:b2:frontier-orbitals:diels-alder}{dienofiel} tegenover zijn uiteinden.
  \item Teken de twee nieuwe $\sigma$-bindingen van de uiteinden van het \omterm{def:b2:frontier-orbitals:diels-alder}{dieen} naar de koolstofatomen van het \omterm{def:b2:frontier-orbitals:diels-alder}{dienofiel},
        en de nieuwe $\pi$-binding tussen C2 en C3 van het \omterm{def:b2:frontier-orbitals:diels-alder}{dieen}.
  \item Houd de substituenten van het \omterm{def:b2:frontier-orbitals:diels-alder}{dienofiel} aan dezelfde kant (cis blijft cis).
  \item Verbind bij asymmetrische partners de grootste coëfficiënten (``ortho''- of ``para''-producten
        voor 1- of 2-gesubstitueerde \omterm{def:b2:frontier-orbitals:diels-alder}{diënen}).
  \item Zet bij een cyclisch \omterm{def:b2:frontier-orbitals:diels-alder}{dieen} de onverzadigde substituent endo.
\end{enumerate}
\end{method}

\begin{inthelab}[Dicyclopentadieen kraken]
Cyclopentadieen wordt verkocht als zijn dimeer en vlak voor gebruik geregenereerd: het dimeer wordt
verhit tot zijn kookpunt (ongeveer \qty{170}{\degreeCelsius}) in een kolf met een
fractioneerkolom, waar het terugsplitst in het monomeer (een \omterm{def:b2:frontier-orbitals:diels-alder}{retro-diels-alderreactie}),
dat bij \qty{41}{\degreeCelsius} overdestilleert en opgevangen wordt in een opvangkolf die in ijs
gekoeld wordt. Het monomeer dimeriseert bij kamertemperatuur binnen enkele uren opnieuw en wordt meteen
gebruikt. Beide verbindingen zijn ontvlambaar en schadelijk; er wordt in een zuurkast gewerkt.
\end{inthelab}

\begin{safety}{\ghs[5mm]{GHS02}\,\ghs[5mm]{GHS06}\,\ghs[5mm]{GHS08}}
Cyclopentadieen en zijn dimeer zijn zeer ontvlambaar en giftig bij inademing; maleïnezuuranhydride
is corrosief en sensibiliseert de luchtwegen. Ze worden in een zuurkast gehanteerd,
met handschoenen en oogbescherming, ver van vlammen.
% ledger: ghs:cyclopentadiene, ghs:dicyclopentadiene, ghs:maleic-anhydride, bp:cyclopentadiene, bp:dicyclopentadiene
\end{safety}

\section{Oefeningen}

\begin{exercise}[$\star$]\label{exo:b2:frontier-orbitals:1}
Geef de \omterm{def:b2:frontier-orbitals:frontier}{HOMO} en de \omterm{def:b2:frontier-orbitals:frontier}{LUMO} (energieën in hückeleenheden) van etheen, van het allylanion
en van butadieen.
\end{exercise}

\begin{exercise}[$\star$]\label{exo:b2:frontier-orbitals:2}
Deel in als hard of zacht: \ce{OH-}, \ce{I-}, \ce{H+}, een thiolaat \ce{RS-}, het koolstofatoom van
joodmethaan, \ce{F-}.
\end{exercise}

\begin{exercise}[$\star$]\label{exo:b2:frontier-orbitals:3}
Welke van deze \omterm{def:b2:frontier-orbitals:diels-alder}{diënen} kunnen deelnemen aan een \omterm{def:b2:frontier-orbitals:diels-alder}{diels-alderreactie}: buta-1,3-dieen,
cyclopentadieen, penta-1,4-dieen, (2\textit{Z},4\textit{Z})-hexa-2,4-dieen in zijn
s-cisvorm?
\end{exercise}

\begin{exercise}[$\star$]\label{exo:b2:frontier-orbitals:4}
Teken het product van butadieen met propeennitril \ce{CH2=CH-C#N}.
\end{exercise}

\begin{exercise}[$\star\star$]\label{exo:b2:frontier-orbitals:5}
Twee orbitalen bij $-10$ en \qty{-2}{eV} wisselwerken met $\beta = \qty{-1.0}{eV}$ (gegevens voor de oefening).
Bereken de storingsstabilisatie van twee elektronen, en daarna de exacte.
\end{exercise}

\begin{exercise}[$\star\star$]\label{exo:b2:frontier-orbitals:6}
Het cyanide-ion reageert met joodmethaan tot ethaannitril \ce{CH3-C#N}, maar met een
carbokation in een sterk ioniserend oplosmiddel geeft het wat isonitril \ce{R-N#C}. Verklaar.
\end{exercise}

\begin{exercise}[$\star\star$]\label{exo:b2:frontier-orbitals:7}
Voorspel met de coëfficiënten van het hoofdstuk het hoofdproduct van 1-methoxybutadieen
met propenal.
\end{exercise}

\begin{exercise}[$\star\star$]\label{exo:b2:frontier-orbitals:8}
Teken het \omterm{def:b2:frontier-orbitals:endo}{endo-adduct} van cyclopentadieen met methylpropenoaat.
\end{exercise}

\begin{exercise}[$\star\star$]\label{exo:b2:frontier-orbitals:9}
Waarom reageert maleïnezuuranhydride bij kamertemperatuur met cyclopentadieen, terwijl etheen
verhitting onder druk nodig heeft?
\end{exercise}

\begin{exercise}[$\star\star\star$]\label{exo:b2:frontier-orbitals:10}
Cyclopentadieen reageert met dimethylmaleaat (cis-diëster) en met dimethylfumaraat
(trans-diëster). Teken beide producten en zeg hoeveel stereo-isomeren elk geeft.
\end{exercise}

\begin{exercise}[$\star\star\star$]\label{exo:b2:frontier-orbitals:11}
De \omterm{def:b2:frontier-orbitals:diels-alder}{diels-alderreactie} heeft $\Delta_r H^\circ < 0$ en $\Delta_r S^\circ < 0$. Verklaar de tekens
en waarom de omgekeerde reactie bij hoge temperatuur gunstig wordt.
\end{exercise}

\begin{exercise}[$\star\star\star$]\label{exo:b2:frontier-orbitals:12}
In hydrazine \ce{H2N-NH2} draagt elk stikstofatoom een vrij elektronenpaar. Leg met de wisselwerking met vier
elektronen uit waarom het molecuul de conformatie vermijdt waarin de twee vrije
elektronenparen parallel staan.
\end{exercise}

\section{Probleem: cyclopentadieen kraken}

\begin{problem}[{Weekendprobleem --- de dimerisatie van cyclopentadieen als
diels-alderreactie, de thermodynamica van het kraken, de grensorbitalen van
cyclopentadieen en maleïnezuuranhydride, en de stereochemie van hun adduct}]
\label{pb:b2:frontier-orbitals:1}
Gegevens: kookpunten, cyclopentadieen \qty{41}{\degreeCelsius}, dicyclopentadieen ongeveer
\qty{170}{\degreeCelsius}. Energieën volgens het hückelmodel (een model): cyclopentadieen, behandeld als
een butadieeneenheid, \omterm{def:b2:frontier-orbitals:frontier}{HOMO} $\alpha + 0.62\beta$ en \omterm{def:b2:frontier-orbitals:frontier}{LUMO} $\alpha - 0.62\beta$; maleïnezuuranhydride, \omterm{def:b2:frontier-orbitals:frontier}{LUMO}
ongeveer $\alpha - 0.35\beta$.
% ledger: bp:cyclopentadiene, bp:dicyclopentadiene

\medskip
\noindent\textbf{Deel I --- Het dimeer.}
\begin{enumerate}
  \item Teken cyclopentadieen en toon aan dat zijn dieeneenheid in de s-cisvorm vastgehouden wordt.
  \item Bij de dimerisatie is het ene molecuul het \omterm{def:b2:frontier-orbitals:diels-alder}{dieen}. Wat is het andere?
  \item Teken het dimeer en nummer de nieuwe bindingen.
  \item Welk adduct, endo of exo, ontstaat bij kamertemperatuur, en waarom?
  \item Is de dimerisatie stereospecifiek? Verklaar met het mechanisme.
  \item Waarom heeft de dimerisatie geen katalysator nodig?
\end{enumerate}

\medskip
\noindent\textbf{Deel II --- Thermodynamica van het kraken.}
\begin{enumerate}[resume]
  \item Geef het teken van $\Delta_r H^\circ$ van de dimerisatie, uit de gevormde en de
        verbroken bindingen.
  \item Geef het teken van $\Delta_r S^\circ$.
  \item Schrijf $\Delta_r G^\circ(T)$ op en toon aan dat het van teken verandert bij een temperatuur $T_i$.
  \item Waarom is het dimeer stabiel bij kamertemperatuur, en wordt het monomeer warm begunstigd?
  \item Waarom drijft het afdestilleren van het monomeer naarmate het ontstaat, in de kraakopstelling de
        reactie voort?
  \item Waarom moet het monomeer koud bewaard en snel gebruikt worden?
\end{enumerate}

\medskip
\noindent\textbf{Deel III --- \omterm{def:b2:frontier-orbitals:frontier}{Grensorbitalen}.}
\begin{enumerate}[resume]
  \item Bereken voor cyclopentadieen dat met zichzelf reageert, de afstand HOMO--LUMO in eenheden van
        $|\beta|$.
  \item Bereken voor cyclopentadieen met maleïnezuuranhydride de afstand tussen de \omterm{def:b2:frontier-orbitals:frontier}{HOMO} van het
        \omterm{def:b2:frontier-orbitals:diels-alder}{dieen} en de \omterm{def:b2:frontier-orbitals:frontier}{LUMO} van het anhydride.
  \item Welk paar reageert sneller? Verklaar.
  \item Waarom verlagen de twee carbonylgroepen van maleïnezuuranhydride zijn \omterm{def:b2:frontier-orbitals:frontier}{LUMO}?
  \item Welke orbitaal van maleïnezuuranhydride geeft de secundaire overlap van de
        endonadering?
\end{enumerate}

\medskip
\noindent\textbf{Deel IV --- Het adduct met maleïnezuuranhydride.}
\begin{enumerate}[resume]
  \item Teken het \omterm{def:b2:frontier-orbitals:endo}{endo-adduct}.
  \item Markeer zijn stereocentra.
  \item Is het adduct chiraal? (Zoek een spiegelvlak.)
  \item Waarom komen de twee waterstofatomen van het vroegere \omterm{def:b2:frontier-orbitals:diels-alder}{dienofiel} aan dezelfde kant van de
        ring terecht?
  \item Hoe zou het \omterm{def:b2:frontier-orbitals:endo}{exo-adduct} eruitzien?
  \item Kunnen het endo- en het \omterm{def:b2:frontier-orbitals:endo}{exo-adduct} bij kamertemperatuur in elkaar overgaan? Waarom?
  \item Hoe zou het anhydride met water reageren, en welk product zou er ontstaan?
  \item Geef het aantal stereocentra van het \omterm{def:b2:frontier-orbitals:endo}{endo-adduct} van cyclopentadieen en
        maleïnezuuranhydride.
\end{enumerate}
\end{problem}
